\subsection{Exact section coordinates with ramification allowed at three and seven}

\begin{statement}{B.33}{The entire finite-set section space and its local fibres}{Proved}
Keep the curve, rational tangent, local factors, and actual
quotient $\Gamma_S$ of B.29--B.30.  Let $Q=N/J$, where
$J$ is any closed normal Galois-stable subgroup; the first
classification below does not require $J\subseteq N'$.
Let $\operatorname{Sect}_S(Q)$ denote the actual sections
of $Q\rtimes G_{\mathbb Q}$ whose whole geometric cocycle
factors through $G_{\mathbb Q,\{3,7,\infty\}}$.
Choose the local generators $t,F$ of $L_1$ in (B.29.2),
put $\Phi=\rho_Q(F)$, and set
\[
 \mathcal Z_Q=
 \{(u,v)\in Q^2:\ v\Phi(u)v^{-1}=u^7\},\qquad
 Q^c=\{w\in Q:\rho_Q(c)(w)=w\}.                              \tag{B.33.1}
\]
There is an exact bijection
\[
        \operatorname{Sect}_S(Q)/Q
             \simeq Q^c\backslash\mathcal Z_Q,                \tag{B.33.2}
\]
where the action on the right is
\[
       w\cdot(u,v)=
       \bigl(wuw^{-1},\;wv\Phi(w)^{-1}\bigr).
                                                               \tag{B.33.3}
\]
The same construction gives a bijection between actual
cocycles normalized by $z(c)=1$ and actual continuous
cocycles $L_1\to Q$; $u=z(t)$ and $v=z(F)$ are their
two explicit local coordinates.

For a local conjugacy class $\xi$, choose a representative
and let $S_\xi\subseteq Q$ be its stabilizer under geometric
conjugation, equivalently the fixed subgroup for its
twisted local action.  The fibre of localization from
(B.33.2) to this local class is
\[
                         Q^c\backslash Q/S_\xi.              \tag{B.33.4}
\]
For the full free pro-three group $Q=N$, every such
$S_\xi$ is trivial.  Thus each local-class fibre is
$N^c\backslash N$, after choosing its representative.
This last assertion is not extended to arbitrary quotients
whose lower central graded groups may have torsion.

All members of the non-rational-cusp family of B.30,
including its zero parameter, fail local geometricity
at three.  The same is true in every quotient with
$J\subseteq N'$.
\end{statement}
\begin{proof}
The actual action over $K$ has pro-three image by B.30.
For any cocycle under consideration its graph on
$G_{K,S_K}$ has target $Q\rtimes\rho_Q(G_{K,S_K})$,
which is pro-three.  Hence its cocycle and action factor
through the same actual quotient $\Gamma_S$.
This is the graph factorization of B.27 applied with the
enlarged finite set; it concerns the whole cocycle.

Write $b=z(c)$.  The cocycle relation gives
$\rho_Q(c)(b)=b^{-1}$.  The square root $b^{1/2}$ exists
uniquely in the procyclic pro-three group generated by
$b$, because $2$ is a unit in $\mathbb Z_3$.
Conjugating the cocycle by $b^{-1/2}$ on the left gives
value one at $c$.  Two normalized cocycles are
geometrically conjugate by $w\in Q$ precisely when
that conjugation has $w\rho_Q(c)(w)^{-1}=1$, namely
$w\in Q^c$.

For a normalized cocycle, restriction to $L_2$ is forced
to be the complex conjugate of its restriction to $L_1$.
Conversely any actual cocycle on $L_1$ extends by this
formula and the free-product isomorphism B.29 to a unique
normalized cocycle on $\Gamma_S$.  Pulling its geometric
component back to $G_{\mathbb Q}$ as in (B.30.5) gives
an actual global section.  Thus no condition of local
geometricity has been inserted into the classification.

Local inertia acts trivially on $Q$.  Evaluating the
actual graph homomorphism on the tame relation gives
\[
 (v,F)(u,t)(v,F)^{-1}
        =(v\Phi(u)v^{-1},t^7),\qquad
                         (u,t)^7=(u^7,t^7).
\]
Consequently (B.33.1) is necessary.
It is also sufficient: the pro-three presentation
(B.29.2) gives a continuous homomorphism
$L_1\to Q\rtimes L_1$ with these two generator images,
and its second projection is the identity.
This verifies continuity and the direction from the
displayed equation to a full local cocycle.
Conjugation by $(w,1)$ gives exactly (B.33.3).
The orbit of a fixed local representative is $Q/S_\xi$,
and then quotienting by $Q^c$ gives (B.33.4).

For $Q=N$ the stabilizer is trivial already under the
twisted action of the single arithmetic Frobenius $F$.
Indeed the closed lower central quotients of a free
pro-three group have torsion-free $\mathbb Z_3$ graded
pieces, with the integral Hall Lie basis.
They can be obtained by collecting Hall commutators
in the free nilpotent quotients and allowing their
integer exponents to converge in $\mathbb Z_3$;
uniqueness follows from their leading Magnus terms.
The first graded piece is the extension of the compact
Tate module by the rank-eight boundary module in
(B.30.6).  The eigenvalues of $F$ on its compact
part have complex absolute value $\sqrt7$, and on
its boundary part have absolute value $7$.
At seven the compact characteristic polynomial is
$X^2+X+7$: the exact point count in
\path{c_seven_free_product.py} is $\#X_3(\mathbb F_7)=9$.
The boundary action is seven times the cusp permutation.
Thus this assertion can also be checked directly here,
without a general weight theorem.
Each degree-$r$ Lie eigenvalue is a product of $r$
first-degree eigenvalues, and hence has complex absolute
value at least $7^{r/2}>1$.
In particular it is not one.
The inner twist contributed by the local cocycle acts
trivially on every lower central graded piece.
Every fixed vector on such a graded piece is therefore
zero: tensoring with $\mathbb Q_3$ has no eigenvalue one,
and the integral piece is torsion-free.
If an element of $N$ were fixed by the twisted Frobenius,
induction would put it in every closed lower central
subgroup.  Their intersection is trivial, since every
finite three-group quotient is nilpotent and these
quotients separate the elements of $N$.
This proves $S_\xi=1$ for the full group.

It remains to verify the stated local failure for B.30.
Write $E=X_3$ with origin $O$ and $T=T_3E$.
Every $s_a$ has the same compact cocycle with values in
$T$: the twists $a\kappa_3(7)$ are peripheral and vanish
on $T$, so their actual compact graph maps agree on both
free factors, and hence globally by uniqueness.
Call its class $h$.  It is unramified outside three,
since its seven-adic inertia is also peripheral.
Its localization at seven is the Kummer class of $P-O$.

The nine coordinate cusps are the nine three-torsion
points relative to $O$.  To see their order directly,
the tangent line at $[0:-\zeta^j:1]$ is
$\zeta^{2j}Y+Z=0$, and its intersection with the cubic
is given by $X^3=0$.  The other coordinate cases are
symmetric.  The ratio of this tangent line to the
tangent line at $O$ therefore has divisor $3P-3O$.
The nine points are distinct and reduce distinctly at
seven.  Exact counting gives nine points in all, so
$E(\mathbb F_7)=E[3]$.
The kernel of good reduction on $E(\mathbb Q_7)$ is
pro-seven.  Its three-adic completion is consequently
$E(\mathbb F_7)$, and the local $T$-Kummer map is
injective on these nine torsion points.

By A.10, $E(\mathbb Q)$ consists of the three rational
coordinate cusps; by A.14 its Tate--Shafarevich group
is finite.  The three-primary Tate Selmer group is
therefore $E(\mathbb Q)^{\wedge}_3$.
If $h$ were locally Kummer at three, it would lie in
that Selmer group: at every other finite prime good
reduction identifies unramified $T$-classes with the
local Kummer image, as in B.27, and the real-place
cohomology vanishes for this odd-primary coefficient.
Its seven-adic class would then be the Kummer image of
one of the three rational cusps.  This contradicts
the preceding injectivity and $P\notin E(\mathbb Q)$.
Hence $h$ is not locally Kummer at three.
Every local point or cusp section has a compact Kummer
class, so none of the $s_a$ is locally geometric there.
For $J\subseteq N'$ the full compact projection is
retained, and the same obstruction applies.
\end{proof}

\begin{statement}{B.34}{The rational local colour can also be noncuspidal globally}{Refuted}
Even after requiring the local cusp at seven to be the
rational base cusp $O$, the rational-cusp implication
of task B.10(b) is false for full $N$ and for every
closed normal stable quotient $N/J$ with $J\subseteq N'$.

Here is a second complete family of counterexamples.
Let $T=T_3E$ and choose $h\in T$ with $c(h)\ne h$.
Write $\operatorname{pr}_E:N\to T$ for compact
abelianization.  Use the entire nonempty fibre
\[
                  \mathcal N_h=\operatorname{pr}_E^{-1}(h)
\]
as the allowed set of lifts: every $n\in\mathcal N_h$
gives the counterexample below.
For an entirely specified lattice choice one may use
$h=(0,1)^t$ in the basis of B.28, where
\[
               C_E=\begin{pmatrix}1&-1\\0&-1\end{pmatrix}.
                                                               \tag{B.34.1}
\]
Alternatively choose a primitive vector in the negative
eigenspace of $c$; the Weil pairing and $2\in\mathbb Z_3^\times$
give an integral decomposition of $T$ into two rank-one
eigenspaces.  On $L_1$ conjugate the local base-cusp
cocycle $a\kappa_3(7)$ by $n$ on the left.  On $L_2$
take its complex conjugate, and set $z(c)=1$.
B.29 gives the resulting actual full section, denoted
$r_{a,n}$, for every $a\in\mathbb Z_3$.
All these sections are unramified outside $\{3,7\}$;
for $a\ne0$ they are ramified at seven with local colour
$O$.  Their compact class has infinite order, and they
are not globally cuspidal.  Using A.10 and A.14 they
also all fail local geometricity at three.
\end{statement}
\begin{proof}
The construction is the literal one in B.30 with a
different prescribed local graph, so its continuity,
factorization, and descent to an actual $G_{\mathbb Q}$
section are already proved.  The compact projection
$N\to T$ is surjective by the homology sequence, so a
lift $n$ of the chosen $h$ exists.
The unit tangent at $O$ is our original splitting.
Its cusp twist has zero compact cocycle.  Left
conjugation by $n$ therefore changes the compact
cocycle on $L_1$ to
\[
        k(F)=(1-\Phi_E)h,\qquad k(t)=0,
                    \Phi_E=\rho_T(F).                        \tag{B.34.2}
\]
The descended cocycle has $k(c)=0$.
It is independent of $a$ and unramified outside three.

Suppose that its class became zero in
$H^1(G_{\mathbb Q},T\otimes\mathbb Q_3)$.
Then for some $m\in T\otimes\mathbb Q_3$ one would have
$k(g)=m-gm$ for every $g$.
The equation at $c$ gives $m=c(m)$, whereas the equation
at $F$ gives
\[
                       (1-\Phi_E)m=(1-\Phi_E)h.
\]
The characteristic polynomial $X^2+X+7$ above shows
$\det(1-\Phi_E)=9\ne0$.
Consequently $m=h$, contradicting $c(h)\ne h$.
The compact class thus has nonzero rational image and
infinite order in $H^1(G_{\mathbb Q},T)$.

Local conjugation has not changed the local cusp
class, and its abelian inertia is $ae_O$.
Every global rational cusp of the cubic has compact
Kummer class killed by three, by the tangent-line
calculation in B.33.  Its rational compact cohomology
class is therefore zero.  The nonzero class just
proved excludes every global cusp packet, for every
$a$, including $a=0$.
For $a\ne0$ the nonzero inertia verifies the required
ramification at seven.

Finally the compact cocycle is unramified outside
three.  If it were locally Kummer at three, good
reduction at all other finite primes would put it
in the Tate Selmer group.  A.10 and A.14 make that
group finite, contradicting the infinite order.
Thus no local point or cusp section can induce
$r_{a,n}$ at three.
All these arguments survive projection through
$J\subseteq N'$, which retains $T$ and $ae_O$.
No local packet classification from an earlier AI
return is assumed: the local sections here are
constructed tangential sections themselves.
\end{proof}

\paragraph{The precise remaining multi-prime problem.}
B.33 classifies the finite set $\{3,7\}$ only.
It does not solve task B.10(a) or the full pro-three
form of B.11.  The failure of the most direct extension
can be quantified.  If $m$ primes of $K$ away from
$(\lambda)$ are allowed, Kummer theory gives
$\dim_{\mathbb F_3}H^1=m+2$, because the localized
class group is still zero.  A free product of all
$m$ corresponding local tame pro-three groups would
have $H^1$ of dimension $2m$.
These dimensions agree only for $m=2$.
Adjoining a new allowed $K$-prime adds one global
generator, whereas adjoining its entire local factor
freely would add two.  Thus an additional Frobenius
word is constrained; the independent local
prescription of B.29 cannot simply be iterated.
Constructing a second colour, or infinitely many
nonzero inertia values tending to one, still requires
solving those actual Frobenius constraints in the
whole free pro-three group.
